\section{Volume and continuum limits}\label{sec:continuum-limits}

Gaussian domination gives compactness of the curvature laws at all
cutoffs. Diagonal reflection positivity uses a further identification:
at a fixed spacing, a volume limit is the step-field image of an
infinite lattice law. This permits the nine-mirror theorem to be
applied before the spacing is removed.

\subsection{The test map and its bounds}

Write $\eta=L^{-N}$, $R=8hL^T$, and $n=R/\eta$, as in
\eqref{eq:cutoff-scales}. A plaquette of the unit covering lattice is
$\sigma=(x;\mu\nu)$, with $\mu<\nu$ and centre
$y_\sigma=x+(e_\mu+e_\nu)/2$. Its averaging cube is
\[
 Q_\eta(\sigma)=\eta\bigl(y_\sigma+[-1/2,1/2)^3\bigr).
\]
For each fixed orientation these cubes partition $\R^3$.
If $u$ is a real lattice curvature, its physical step field is
$\eta^{-3/2}u(\sigma)$ on $Q_\eta(\sigma)$ in component $\mu\nu$.
For a Schwartz two-form $f$, define
\begin{align}
 f_{[\eta]}(\sigma)
   &=\eta^{-3/2}\int_{Q_\eta(\sigma)}f_{\mu\nu}(y)\,\dd y,
                                                   \label{eq:limit-test-cover}\\
 f^R(y)&=\sum_{k\in\Z^3}f(y+Rk),\notag\\
 f_\eta([\sigma])
   &=\eta^{-3}\int_{Q_\eta(\sigma)}f^R_{\mu\nu}(y)\,\dd y.
                                                   \label{eq:limit-test-torus}
\end{align}
Here $[\sigma]$ is the plaquette class modulo $n\Z^3$.
The torus cochain norm is $|v|_\eta^2=\eta^3\sum_{[\sigma]}|v([\sigma])|^2$.
The covering cochain norm $\|\cdot\|$ is the unweighted $\ell^2$ norm.
For the periodic extension $\widetilde u$ of $u$,
\begin{equation}\label{eq:limit-step-pairing}
 F^\sharp(f)=\eta^{3/2}\sum_{[\sigma]}f_\eta([\sigma])u(\sigma)
            =\sum_\sigma f_{[\eta]}(\sigma)\widetilde u(\sigma).
\end{equation}
The normalization in the first equality combines the cell volume
$\eta^3$ with the physical field factor $\eta^{-3/2}$.

The estimates needed below are
\begin{align}
 \|f_{[\eta]}\|&\le\|f\|_2, \label{eq:limit-cover-norm}\\
 \sum_\sigma|f_{[\eta]}(\sigma)|(1+|x(\sigma)|_\infty)
 &\le\eta^{-3/2}\int_{\R^3}\sum_{\mu<\nu}|f_{\mu\nu}(y)|
                         (2+|y|_\infty/\eta)\,\dd y, \label{eq:limit-cover-sum}\\
 |f_\eta|_\eta&\le\|f^R\|_{L^2(\R^3/R\Z^3)}\le p_R(f),\label{eq:limit-torus-norm}
\end{align}
where, with $Q_R=[-R/2,R/2)^3$,
\begin{equation}\label{eq:limit-periodic-seminorm}
 p_R(f)=\|f\|_2+
        26R^{-2}\||y|^2f\|_{L^2(\R^3\setminus Q_R)}.
\end{equation}
One also has $|f_\eta|_\eta\le\|f\|_2$ if the support of $f$ lies
in a translate of an open cube of side $R$. For $R\ge1$,
$p_R(f)\le p_1(f)$, and $p_R(f)\to\|f\|_2$ as $R\to\infty$.
The seminorm $p_1$ is continuous on Schwartz space.

The first bound follows from Cauchy--Schwarz on each averaging
cube. On that cube $1+|x(\sigma)|_\infty\le2+|y|_\infty/\eta$,
which gives the second. Absolute convergence permits unfolding:
\[
 \eta^{3/2}f_\eta([\sigma])
       =\sum_{k\in\Z^3}f_{[\eta]}(\sigma+nk).
\]
This proves \eqref{eq:limit-step-pairing}. Applying the cube
Cauchy--Schwarz inequality to $f^R$ gives the first inequality in
\eqref{eq:limit-torus-norm}. If the support lies in one open
$R$-cube, its distinct $R\Z^3$ translates are disjoint, so
$\|f^R\|_{L^2(\R^3/R\Z^3)}=\|f\|_2$.

For the remaining estimate put $a_k=\|f\|_{L^2(Q_R+Rk)}$.
Minkowski's inequality bounds the periodized norm by $\sum_k a_k$.
On $Q_R+Rk$, $k\ne0$, one has $|y|\ge R|k|_\infty/2$; hence
\[
 \sum_{k\ne0}a_k
 \le4R^{-2}\left(\sum_{k\ne0}|k|_\infty^{-4}\right)^{1/2}
                   \||y|^2f\|_{L^2(\R^3\setminus Q_R)}.
\]
The number of points with $|k|_\infty=j$ is $24j^2+2$, so
the sum in parentheses is $4\pi^2+\pi^4/45<41.65$ and its
square root times four is less than $26$.
This proves \eqref{eq:limit-torus-norm} and
\eqref{eq:limit-periodic-seminorm}. The comparison with $p_1$
and the limit of $p_R$ follow by inclusion of the cubes and
integrability of $|y|^4|f(y)|^2$.

For the Hermite two-forms of \cref{sec:probability},
\begin{equation}\label{eq:limit-Hermite-norm}
 p_1(h_\kappa)\le31\lambda_\kappa.
\end{equation}
Here is the weighted estimate behind this bound. Let
$\mathsf M=1+|y|^2$, $P_j=-i\partial_j$, and
$\mathsf A=\mathsf M-\Delta$. Since
$[P_j,[P_j,\mathsf M]]=-2$,
\[
 \|\mathsf A v\|_2^2
 =\|\mathsf Mv\|_2^2+\|\Delta v\|_2^2
      +2\sum_j\langle P_jv,\mathsf MP_jv\rangle-6\|v\|_2^2
 \ge\|\mathsf Mv\|_2^2-6\|v\|_2^2.
\]
The Hermite eigenvalues satisfy $\lambda_\kappa\ge4$, and therefore
$\|\mathsf Mv\|_2^2\le(11/8)\|\mathsf Av\|_2^2$.
Taking $v=h_\kappa$ gives
$p_1(h_\kappa)\le1+26\sqrt{11/8}\lambda_\kappa
\le31\lambda_\kappa$.

The test map respects the affine cubic symmetries. For a signed
permutation and lattice translation $g$ of the unit lattice, let
$\mathfrak g(y)=\eta g(y/\eta)$. With the exterior-square action
on two-forms and the cochain pull-back,
\begin{equation}\label{eq:limit-test-equivariance}
 (\mathfrak g\cdot f)_{[\eta]}=(g^{-1})^*f_{[\eta]},
 \qquad
 (\mathfrak g\cdot f)_\eta=(g^{-1})^*f_\eta.
\end{equation}
Indeed, $g$ sends centres to centres and the averaging cubes to
averaging cubes up to null boundaries. Its orientation sign is
the sign of $\wedge^2\vec g$ on the corresponding plaquette.
Changing variables gives the first formula; periodization and
the bijection $k\mapsto\vec gk$ give the second.

\subsection{Compactness and convergence of moments}

Let $P_{N,T}$ be the law of $F^\sharp$ for the all-flux chamber
measure. The mass assignment may be any real sequence independent of
$T$. Proposition~\extref{5.1}{ChamberLattice} and
\eqref{eq:limit-step-pairing} give
\begin{equation}\label{eq:limit-cutoff-G}
 \E_{N,T}e^{F^\sharp(f)}
       \le e^{|f_\eta|_\eta^2/2}\le e^{p_R(f)^2/2}
       \qquad(f\text{ real}).
\end{equation}
The random-linear-form proof of Proposition~\extref{2.1}{ContinuumLimits}
applies here with the Hilbert seminorm $f\mapsto|f_\eta|_\eta$.
In particular its variance bound and \eqref{eq:limit-Hermite-norm}
give
\begin{equation}\label{eq:limit-uniform-tightness}
 \E_{N,T}\|F^\sharp\|_{-3}^2
 \le\sum_\kappa\lambda_\kappa^{-6}p_1(h_\kappa)^2
 \le961\sum_\kappa\lambda_\kappa^{-4}<61.
\end{equation}
Tonelli's theorem justifies the first sum. Every finite lattice
cochain has a step field in $H_{-3}$: Cauchy--Schwarz in
\eqref{eq:limit-step-pairing} bounds its $\kappa$-th Hermite
coordinate by $31\lambda_\kappa\|u\|_{\ell^2(C(n)^{(2)})}$,
with the unweighted norm on the unit-lattice torus plaquettes.
The step map from this finite-dimensional cochain space to $H_{-3}$
is consequently continuous.

The compact embedding recalled in \cref{sec:probability}
now applies. The laws $P_{N,T}$ are uniformly tight on $H_{-4}$,
and their fixed-spacing volume limits are uniformly tight there.
Thus the sets $\mathcal V_N$ and the set of volume-first limits
are nonempty.

For a joint sequence $N_k,T_k\to\infty$ suppose
$P_{N_k,T_k}\to\mu$ weakly on $H_{-4}$. Truncating $e^{F(f)}$
above by a constant, passing to the weak limit, and then removing
the truncation gives
\begin{equation}\label{eq:limit-continuum-G}
 \E_\mu e^{F(f)}\le e^{\|f\|_2^2/2}.
\end{equation}
The limiting constant follows from $p_{R_k}(f)\to\|f\|_2$.
For a fixed finite family of complex tests $f_1,\ldots,f_s$, define
\[
 Z_k(z)=\E_{N_k,T_k}\exp\left(\sum_a z_aF^\sharp(f_a)\right),
 \qquad
 Z(z)=\E_\mu\exp\left(\sum_a z_aF(f_a)\right).
\]
The evaluations $F\mapsto(F(f_a))_a$ are continuous on $H_{-4}$.
Moreover,
\[
 \E_{N_k,T_k}\left|\exp\left(\sum_a z_aF^\sharp(f_a)\right)\right|^2
 \le\exp\left(2p_1\left(\operatorname{Re}\sum_a z_af_a\right)^2\right).
\]
This bound is uniform in $k$ and on compact sets of $z$.
Uniform integrability therefore gives pointwise convergence
$Z_k\to Z$. The same exponential bounds give local uniform
bounds on the entire functions $Z_k$. Cauchy's estimates make them
equi-Lipschitz on compact sets, so their convergence is locally
uniform. Cauchy's formula then passes every derivative to the limit.
In particular all joint moments converge.

On each compact subset of
$\{\sum_a|z_a|\|f_a\|_2<2/5\}$, the test-map bounds give
$\sum_a|z_a||(f_a)_{\eta_k}|_{\eta_k}<2/5$ for all large $k$.
The zero-free estimate \exteq{17}{ContinuumLimits}, with this
seminorm, bounds $|Z_k-1|$ by $1/2$. The principal logarithms
therefore converge locally uniformly near zero, and their
derivatives give convergence of every joint cumulant.

\subsection{The prescribed parent and its components}

The finite-periodic moment passage just proved applies to a
diagonal realization of any prescribed volume-first limit.
The fixed-spacing recovery in \cite[Proposition~7.3]{ContinuumLimits}
identifies each intermediate volume law with the interpolation
of an infinite-lattice law. Its nine reflection forms then pass
to the continuum limit by \cite[Corollary~7.4]{ContinuumLimits}.
We record the resulting input to the construction.

\begin{proposition}\label{prop:limit-nine-mirrors}
Let $\mu$ be any volume-first limit for a real, volume-independent
mass assignment. Then $\mu$ is carried by closed tempered two-forms,
is translation invariant and cubic invariant, and satisfies
\eqref{eq:limit-continuum-G}. It is invariant and reflection
positive at each plane in the nine families
\[
 x_j=c\quad(1\le j\le3),\qquad
 x_j\pm x_k=c\quad(1\le j<k\le3),\qquad c\in\R.
\]
There is a sequence of finite laws $P_{N_i,T_i}\to\mu$ with
$N_i,T_i\to\infty$ along which every fixed finite-family moment
and cumulant converges. In the translation-ergodic decomposition
of $\mu$, a single full-measure set of components has the same
closedness, Gaussian bound, cubic symmetries and nine reflection
forms. Its component laws have clustering for all $L^2$ observables
in all directions.
\end{proposition}

\begin{proof}
The parent-law conclusions and the weak finite-law diagonal are
\cite[Corollary~7.4]{ContinuumLimits}; its hypotheses hold by
the all-flux lattice theorem of \cite[Theorem~A]{ChamberLattice}
and \eqref{eq:limit-cutoff-G}. The preceding moment passage
applies to that diagonal. The component conclusions are
\cite[Theorem~A]{ErgodicComponents}, with closedness taken as the
specified event of full parent measure. That theorem gives all
real Schwartz Gaussian bounds simultaneously and the full
translation action on the same component set.
\end{proof}

